% Local construction proofs in the fixed CCHM arXiv-v1 calculus.
% These are not a normalization, canonicity or global model proof.
\begin{annexitem}{annex:p07.cchm2016.singleton}
 {thm:p07-cchm2016-singleton}{All faces of the singleton contraction}
\begin{proof}
Work in \(\Gamma,b:A,p:\mathsf{Path}\,A\,a\,b,i:\mathbb I\).
The first component \(p\,i\) has type \(A\).  In the extension by
\(j:\mathbb I\), \(p(i\wedge j):A\), with \(j=0\) value \(a\)
and \(j=1\) value \(p\,i\).  Path introduction therefore gives
\(\langle j\rangle p(i\wedge j):
\mathsf{Path}\,A\,a\,(p\,i)\), the required dependent second component.
Pair introduction now gives an element of the singleton along \(i\).
At \(i=0\) the pair is
\((a,\langle j\rangle a)=(a,1_a)\).  At \(i=1\) it is
\((b,\langle j\rangle p\,j)=(b,p)\), using path eta.
Name abstraction in \(i\), followed by ordinary abstraction in \(b,p\),
gives the claimed contraction function.  The construction is by
capture-avoiding name binding and the displayed lattice equations;
it uses neither an identity eliminator nor composition.
\end{proof}
\end{annexitem}

\begin{annexitem}{annex:p07.cchm2016.filling}
 {thm:p07-cchm2016-filling}{The complete cap, overlap and endpoint check}
\begin{proof}
Choose \(j\) fresh for \(\Gamma,i\).  The family
\(A(i/i\wedge j)\) has bottom \(A(i/0)\) and top \(A\).
On the side \(\varphi\), the substituted tube is well typed in this
family.  On the side \(i=0\), the family is constantly \(A(i/0)\),
so \(a_0\) has the required type.  Their overlap is
\(\varphi\wedge(i=0)\); there the first branch is \(u(i/0)\),
equal to \(a_0\) by the original cap compatibility.  Thus the two
branches define a tube on \(\varphi\vee(i=0)\).
At \(j=0\) both branches agree with \(a_0\): this is the original cap
condition on \(\varphi\), and reflexivity on \(i=0\).
The composition rule therefore gives a term in the top family \(A\).
Its side equation restricts it to \(u(i/i\wedge1)=u\) on \(\varphi\),
and to \(a_0\) on \(i=0\).  At \(i=1\), its family is \(A(i/j)\),
its tube is \(u(i/j)\) on \(\varphi\), the extra side \(i=0\)
is empty, and its cap is \(a_0\).
Deleting an empty system branch and renaming \(j\) to \(i\) gives
the original composition term.  These operations are typed system
equality, substitution and alpha conversion, not a regularity axiom
for constant transport.  This proves all three judgmental equations.
\end{proof}
\end{annexitem}

\begin{annexitem}{annex:p07.cchm2016.extension}
 {thm:p07-cchm2016-extension}{Extension and the converse with its uniformity quantifier}
\begin{proof}
Let \(c=(c_0,h):\mathsf{isContr}(A)\).
On \(\varphi\), \(h\,u:\mathsf{Path}\,A\,c_0\,u\), so the
tube \(h\,u\,i\) has bottom \(c_0\) and top \(u\).
Composition with cap \(c_0\) is well typed and equals \(u\) on
\(\varphi\).  This proves extension and its strict face equality.
All uses are the specified syntax operations, so the expression is
stable under the substitution equations of the fixed calculus.

Conversely, suppose an extension operation is available for every
face and partial \(A\)-element in every context extension, with strict
restriction and substitution stability.  Set its center to the
extension of the empty partial element.  In the context with \(x:A\)
and a fresh \(i\), form the partial element having value the center
on \(i=0\) and value \(x\) on \(i=1\).  The faces are disjoint and
their join is its declared extent, so this is a valid partial system.
Extend it and abstract \(i\); strict restriction supplies a path from
the center to \(x\).  Abstract \(x\) and pair with the center.
This gives \(\mathsf{isContr}(A)\).
The requirement about all context extensions is essential: an
existential choice of an extension only in the original context
does not give the term with variable \(x\) needed by this argument.
Applying the forward construction to each fiber of \(f\), or the
converse to its uniform partial-fiber extension operation, gives the
corresponding two directions for \(\mathsf{isEquiv}\).
\end{proof}
\end{annexitem}

\begin{annexitem}{annex:p07.cchm2016.preservation}
 {thm:p07-cchm2016-preservation}{The preservation path, including its constant face}
\begin{proof}
The filling \(v:T\) has bottom \(t_0\), top the specified composition,
and restriction \(t\) on \(\varphi\).
The expression \(f\,v:A\) is available throughout \(i,j\); restricting
it to \(\varphi\vee(j=1)\) is legitimate.  At \(i=0\) it equals
\(f(i/0)t_0\), giving the cap compatibility on the whole declared
extent.  The composition is thus well typed in \(A(i/1)\).
At \(j=0\), its extent becomes \(\varphi\), and its side agrees with
\(f\,t\) there, giving the first endpoint.
At \(j=1\), its extent becomes full, so it is
\(f(i/1)v(i/1)\), giving the second endpoint.
Under \(\varphi\), the extent is full for every \(j\), hence the
composition is always \(f(i/1)t(i/1)\).  Its abstraction is therefore
the constant path on that value.  At this face the two endpoints
also equal that same value by composition's side equation.
\end{proof}
\end{annexitem}

\begin{annexitem}{annex:p07.cchm2016.face-decomposition}
 {thm:p07-cchm2016-face-decomposition}{The finite distributive-lattice argument}
\begin{proof}
Distributivity writes a face expression as a finite join of meets of
face generators.  Any meet containing both \(i=0\) and \(i=1\)
is zero and may be deleted.  Divide the remaining meets into those
containing neither \(i\)-face, those containing \(i=0\), and those
containing \(i=1\).  The first class is exactly the image under the
lattice homomorphism sending the two \(i\)-faces to zero.
The other two classes occur, respectively, in
\(\varphi\wedge(i=0)\) and \(\varphi\wedge(i=1)\).
Conversely, each of those three summands is below \(\varphi\):
the two intersections obviously are, and replacing a generator by
zero can only decrease a lattice expression.  This proves equality.

Write \(\delta=\forall i\,\varphi\).  It is independent of \(i\)
and is below \(\varphi\).  If an independent \(\theta\leq\varphi\)
is given, apply the monotone homomorphism \(\forall i\) to obtain
\(\theta=\forall i\,\theta\leq\delta\).  Hence \(\delta\) is largest.
Finally inspect the same meets after intersection with \(i=\varepsilon\).
A meet containing the opposite face becomes zero.  A surviving meet
has its \(i=\varepsilon\) factor removed by substitution and its
other factors unchanged; its intersection before substitution is
below that substituted meet.  Taking joins proves
\(\varphi\wedge(i=\varepsilon)\leq\varphi(i/\varepsilon)\).
No complement operation or excluded-middle law for faces was used.
\end{proof}
\end{annexitem}

\begin{annexitem}{annex:p07.cchm2016.glue-composition}
 {thm:p07-cchm2016-glue-composition}{All caps and overlaps of the corrected Glue operation}
\begin{proof}
We use the notation of the body.  Unglue sends the tube and cap to
\(a:A\) under \(\psi,i\) and \(a_0:A(0)\), agreeing under \(\psi\)
at \(i=0\).  Thus \(a'_1:A(1)[\psi\mapsto a(1)]\).
Under \(\delta\), the gluing face holds throughout the direction.
There \(G=T\), the tube and cap are \(T\)-terms, and
\(a=f.1\,b,\ a_0=f(0).1\,b_0\).  Hence \(t'_1:T(1)\) and the
preservation lemma gives
\(\omega:\mathsf{Path}\,A(1)\,a'_1\,(f(1).1\,t'_1)\).
On \(\delta\wedge\psi\), \(t'_1=b(1)\), \(a'_1=a(1)=f(1).1\,b(1)\),
and \(\omega\) is the constant path on this value.

The final gluing face \(\varphi(1)\) contains \(\delta\).
On its intersection with \(\psi\), \(b(1):T(1)\) and
\(f(1).1\,b(1)=a(1)=a'_1\).  Thus
\((b(1),\langle j\rangle a'_1)\) is a point in the fiber of \(f(1).1\)
over \(a'_1\).  The preceding paragraph checks its overlap with
\((t'_1,\omega)\).  Fiber extension therefore gives a well-typed
\((t_1,\alpha)\), where
\(\alpha:\mathsf{Path}\,A(1)\,a'_1\,(f(1).1\,t_1)\), strictly
extending both prescribed branches.

In the last composition the \(\varphi(1)\)-branch is \(\alpha j\);
the \(\psi\)-branch is \(a(1)\).  On their overlap, extension has made
\(\alpha\) the constant path on \(a'_1=a(1)\), so they agree.
Their bottoms are \(a'_1\), on the first face by the path endpoint
and on the second by \(a'_1=a(1)\).
Composition gives \(a_1:A(1)\), with \(a_1=f(1).1\,t_1\) on
\(\varphi(1)\) and \(a_1=a(1)\) on \(\psi\).
Glue introduction yields \(b_1:G(1)\).
Under \(\psi\) it is
\(\mathsf{glue}[\varphi(1)\mapsto b(1)]
(\mathsf{unglue}\,b(1))=b(1)\) by Glue eta.
Under \(\delta\), the final Glue face is full and \(t_1=t'_1\),
so \(b_1=t'_1\).
If \(\varphi=1\) throughout, \(\delta=1\), which proves the stated
full-face reduction to \(T\)'s composition.
This is a typing and boundary proof within the fixed calculus, not
the separate semantic theorem that every restriction map preserves
the constructed composition structure.
\end{proof}
\end{annexitem}

\begin{annexitem}{annex:p07.cchm2016.unglue}
 {thm:p07-cchm2016-unglue}{The complete partial-fiber extension for unglue}
\begin{proof}
Fix \(a:A\).  In any context extension let a face \(\psi\) carry
\(b:G\) and
\(\alpha:\mathsf{Path}\,A\,a\,(\mathsf{unglue}\,b)\).
On \(\varphi\), \(G=T\) and unglue equals \(f.1\).
Use fiber extension for \(f\) on the partial \(\psi\)-data
\((b,\alpha)\).  It supplies \(t:T\) on \(\varphi\), with \(t=b\)
on \(\varphi\wedge\psi\), and
\(\beta:\mathsf{Path}\,A\,a\,(f.1\,t)\), with
\(\beta=\alpha\) on that intersection.

The tube \([\varphi\mapsto\beta i,\psi\mapsto\alpha i]\) is compatible
on its overlap.  Both branches have bottom \(a\), so let
\(\widetilde a=\mathsf{comp}^{i}A
[\varphi\mapsto\beta i,\psi\mapsto\alpha i]a\).
Its \(\varphi\)-restriction is \(f.1\,t\), allowing
\(\widetilde b=\mathsf{glue}[\varphi\mapsto t]\widetilde a:G\).
Under \(\psi\), \(\widetilde a=\mathsf{unglue}\,b\) and \(t=b\)
where required, hence \(\widetilde b=b\) by Glue eta.
The filling of the same box, abstracted in \(i\), gives
\(\widetilde\alpha:\mathsf{Path}\,A\,a\,\widetilde a\).
By unglue beta its last endpoint is
\(\mathsf{unglue}\,\widetilde b\); by filling's side equation it
equals \(\alpha\) under \(\psi\).
Thus \((\widetilde b,\widetilde\alpha)\) extends the partial element
of the fiber strictly.
These are syntax operations available in every context extension
and stable by their substitution equations.
The converse partial-extension argument of
\ref{thm:p07-cchm2016-extension} now contracts each fiber.
Abstracting \(a\) supplies the required \(\mathsf{isEquiv}\)-term.
This establishes existence of equivalence data for unglue; it does
not assert that the selected data judgmentally restrict to the
original \(f.2\) on \(\varphi\).
\end{proof}
\end{annexitem}

\begin{annexitem}{annex:p07.cchm2016.line-equivalence}
 {thm:p07-cchm2016-line-equivalence}{The original fiber contraction with every boundary checked}
\begin{proof}
Fix \(y:B\).  Use the terms \(f,g,u,v\) of the body and choose fresh
names \(j,k\).  Define the filling
\[
 \theta_0=
 \mathsf{fill}^{i}E
 [(j=0)\mapsto v\,y,(j=1)\mapsto u(g\,y)](g\,y).
\]
Both side branches have bottom \(g\,y\).  They have disjoint faces,
so filling is well typed.  Its \(i=0\) value is \(g\,y\), its
\(j=0\) edge is \(v\,y\), and its \(j=1\) edge is \(u(g\,y)\).
Consequently
\(\gamma=\langle j\rangle\theta_0(i/1):
\mathsf{Path}\,B\,y\,(f(g\,y))\).
The pair \((g\,y,\gamma)\) is a center of the required fiber.
Here and below an endpoint such as \(\theta_0(0)\) means
substitution \(\theta_0(i/0)\); \(\omega(0)\) and \(\omega(1)\)
mean substitution in its remaining \(j\)-direction.

Let \((x,\beta)\) be any point of that fiber, so
\(\beta:\mathsf{Path}\,B\,y\,(f\,x)\).  Fill backwards:
\[
 \theta_1=
 \bigl(\mathsf{fill}^{i}E(i/1-i)
 [(j=0)\mapsto(v\,y)(i/1-i),
  (j=1)\mapsto(u\,x)(i/1-i)](\beta\,j)\bigr)(i/1-i).
\]
Before the final reversal its bottom lies in \(B\).
On \(j=0\) it is \(\beta0=y=v(1)y\); on \(j=1\) it is
\(\beta1=f\,x=u(1)x\).  This checks both cap conditions.
After reversal, \(\theta_1:E\) has \(i=1\) value \(\beta j\),
\(j=0\) edge \(v\,y\), and \(j=1\) edge \(u\,x\).
Set \(\omega=\theta_1(i/0):A\).  Then
\(\omega(j/0)=g\,y\) and \(\omega(j/1)=x\).

Define the following top composition in \(B\), with \(j,k\) free:
\[
 d=\mathsf{comp}^{i}E
 [(k=0)\mapsto\theta_0,\ (k=1)\mapsto\theta_1,\
  (j=0)\mapsto v\,y,\
  (j=1)\mapsto u\,\omega(j/k)]\,
       \omega(j/j\wedge k).
\]
Its cap is an \(A\)-term.  At \(k=0\) it is \(g\,y=\theta_0(0)\);
at \(k=1\) it is \(\omega=\theta_1(0)\);
at \(j=0\) it is \(g\,y=v(0)y\);
at \(j=1\) it is \(\omega(j/k)=u(0)(\omega(j/k))\).
This checks all four cap restrictions.
Of the six pairwise side intersections, \(k=0\) with \(k=1\)
and \(j=0\) with \(j=1\) are empty.  On \(k=0,j=0\) and \(k=1,j=0\),
both selected branches equal \(v\,y\).
On \(k=0,j=1\) they equal \(u(g\,y)\), since \(\omega(0)=g\,y\).
On \(k=1,j=1\) they equal \(u\,x\), since \(\omega(1)=x\).
All overlaps therefore agree, and the box is well typed.

At its top the \(j=0\) face is \(y\), and the \(j=1\) face is
\(f(\omega(j/k))\).  Thus
\(\langle j\rangle d:
\mathsf{Path}\,B\,y\,(f(\omega(j/k)))\) in the context with \(k\).
Pairing it with \(\omega(j/k):A\) gives a point of the fiber along
\(k\).  At \(k=0\) its pair is
\((g\,y,\langle j\rangle\theta_0(i/1))=(g\,y,\gamma)\);
at \(k=1\) it is
\((x,\langle j\rangle\theta_1(i/1))=(x,\beta)\) by path eta.
Hence
\(\langle k\rangle(\omega(j/k),\langle j\rangle d)\)
is the contraction path from the center to \((x,\beta)\).
Ordinary abstraction over the fiber point yields its contraction
function, pairing yields \(\mathsf{isContr}\), and abstraction over
\(y\) yields \(\mathsf{isEquiv}(A,B,f)\).
Every family and endpoint in this construction has been specified;
no assumption that forward and backward transports are inverse was
made.
\end{proof}
\end{annexitem}
